\section{A uniform exact rectangular backend}
\label{sec:matrix}

\subsection{Two fixed shapes suffice}

Fix $\eps=10^{-5}$. The imported characteristic-zero results give, with
any fixed positive slack, square exponent $9/4$ and exponent two for
the shape $N\times N^{1/2}$ by $N^{1/2}\times N$
\cite{uznanski}. We do not assume that an exponent infimum is attained
by one finite algorithm with zero slack.

\begin{lemma}[Uniform endpoint interpolation]\label{lem:rectangle}
For sorted nonnegative logarithmic dimensions $x\ge y\ge z$, the product
of matrices with those three index dimensions can be computed, up to
polynomial bit overhead and fixed rounding factors, with exponent
\begin{equation}\label{eq:rectangle}
 F(x,y,z)=\max\left\{
    (1+\eps/2)(x+y),\quad
    x+(3/4+\eps)y+\tfrac12z\right\}.
\end{equation}
A fixed finite library of algorithms suffices for all shapes occurring in
the contractor. Working storage is a polynomial factor times the largest
input/output array, whose logarithmic size is $x+y$.
\end{lemma}

\begin{proof}
Choose finite bilinear endpoint algorithms with fixed positive slacks at
most $\eps$ in their exponents. First block the largest dimension, paying
$x-y$. If $z\ge y/2$, put
\[
 t_0=2(y-z),\qquad t_1=2z-y.
\]
Both are nonnegative, $t_0+t_1=y$, and $t_0/2+t_1=z$.
Tensoring thin and square algorithms gives exponent at most
\[
 x-y+(2+\eps)t_0+(9/4+\eps)t_1
 =x+(3/4+\eps)y+z/2.
\]
If $z<y/2$, pad the inner dimension to $y/2$. The exponent is at most
$x+y+\eps y\le(1+\eps/2)(x+y)$. This proves the cost bound.

Only two endpoint shapes are used. Choose constant base dimensions
compatible with the rational thin aspect ratio, and round the number of
tensor iterations upward. Each rounding inflates all dimensions by
fixed factors. If necessary allocate a smaller fixed slack to the
endpoint choices to absorb these fixed overheads. No algorithm chosen
from a continuum of input-dependent shapes is required.

Execute the fixed bilinear recursions depth first, processing one child
product at a time and accumulating into output arrays. At every recursive
level, the required linear forms and output arrays are a fixed number
of arrays at that level's dimensions. Their storage sums geometrically
along an active path. Blocking is also sequential. Thus storage is a
constant or logarithmic factor times the largest input/output array.
Permuting the three tensor legs justifies sorting the dimensions.
\end{proof}

For comparison, using only the formal-scope bounds $\omega(\C)\le9/4$
and $\alpha>0.465$ permits the conservative thin ratio $23/50$ and gives
\[
 \max\left\{(1+\eps/2)(x+y),\quad
 x+(85/108+\eps)y+(25/54)z\right\}.
\]
The numerical headline uses the stronger half-width endpoint explicitly
cited above; the weaker expression is recorded to make this dependency
visible.

\subsection{From complex identities to exact modular computation}

The arithmetic model in the source matrix theorems is not itself a
bit-complexity theorem for $R_p$. The following transfer closes that gap.

\begin{lemma}[Fixed algebraic coefficients]\label{lem:descent}
Fix a finite matrix-multiplication bilinear algorithm that exists over
$\C$. An algorithm with the same permitted bilinear support and operation
pattern, and no larger rank, exists with coefficients in a fixed number
field. Outside a fixed finite
set of primes, its identities hold in a fixed-rank finite algebra over the
prime field. These constants and this exceptional set are independent of
the input matrix dimensions obtained by recursively composing the base
algorithm.
\end{lemma}

\begin{proof}
Correctness of the fixed bilinear formula is a finite system of polynomial
equations with rational coefficients in the algorithm's constants: equate
the coefficient of every input monomial to the required multiplication
coefficient. Existence of a complex solution means the ideal is proper.
By the weak Nullstellensatz over $\overline\Q$, it has an algebraic-number
solution. The finitely many coordinates lie in one finite number field
$K$. Fix an integral primitive element $\theta$, a monic integer minimal
polynomial $g$, and an integer $D>0$ such that all formula constants lie
in $\Z[\theta,1/D]$.

For a prime $q\nmid D$, reduce the coefficients in
\[
 A_q=\F_q[T]/(g(T)).
\]
The bilinear identity holds there because it is a polynomial identity.
The polynomial $g$ need not remain irreducible or separable modulo $q$:
the computation is division free after the fixed scalar denominators
have been inverted. The scalar inclusion $\F_q\hookrightarrow A_q$ is
injective, since $g$ is monic of positive degree. The algebra has fixed
dimension, so arithmetic has constant overhead in prime-field operations.
Recursive composition uses this same field and these same constants.
\end{proof}

The lemma is an existence theorem with a uniform algorithm once its finite
base data are fixed. It does not give small base cases or usable practical
crossover sizes. If effectivity of the base choice is desired, algebraic
number coefficients and finite decompositions can be enumerated and their
polynomial identities checked exactly; the existence theorem guarantees
termination before the SAT input is processed. Its cost is a fixed
constant for a fixed exponent slack. No input-dependent search is hidden
inside the asserted exponent.

\begin{proposition}[Transfer to the contractor ring]\label{prop:ring}
The backend of Lemma~\ref{lem:rectangle} applies to the matrices over
$R_p$ used by the contractor with a polynomial factor in $s+\mathcal L$
and the stated input/output storage bound.
\end{proposition}

\begin{proof}
Use the specified free-module basis of $R_p$ over $\Z/2^p\Z$ and lift
each coordinate of each matrix entry to its integer representative in
$[0,2^p)$. For an inner matrix dimension $d$, an ordinary scalar coordinate
product has nonnegative entries bounded by
\[
 d(2^p-1)^2.
\]
Its bit length is $O(p+\log d)$, polynomial in the root parameter. Compute
each such integer matrix product modulo sufficiently many nonexceptional
odd primes so that their product exceeds this bound. For every prime,
Lemma~\ref{lem:descent} supplies the same finite-base algorithm in $A_q$;
read its scalar output coordinates in $\F_q$. Chinese remaindering then
recovers the \emph{integer} coordinate matrix products exactly.

Finally combine the coordinate-pair products with the ring's integer
structure constants and reduce modulo $2^p$. There are only polynomially
many coordinate pairs, and only polynomially many primes of polynomial
bit length are needed. They can be processed sequentially, retaining
partial CRT outputs of polynomial bit length. Thus both the time overhead
and the storage overhead are polynomial.

This argument does not posit a ring homomorphism from $R_p$ to an odd
prime field; no such map is used. It first computes ordinary integer
products of canonical coordinates and only then restores the ring
multiplication. It also never divides by two or by a nonunit of $R_p$.
\end{proof}

The full continuum of shape bounds in~\cite{uznanski} is unnecessary for
this theorem. Using a fixed endpoint library is what supplies uniformity
for the input-dependent logarithmic dimensions of the actual mode lists.
